\labsection{3}{Shell scripting for OS observation and automation}{sec:lab03}
\begin{labproject}{Automate a system report, process query, and log summary}
\textbf{Suggested time:} 90 minutes \hfill
\textbf{Level:} introductory\par
\textbf{Starter files:} \texttt{system\_report.sh}, \texttt{process\_summary.sh}, \texttt{log\_summary.sh}, \texttt{events.log}.\par
\textbf{Evidence to keep:} generated report, process query results, log summary, and one script improvement.

\begin{labmode}
This lab uses Bash for tasks that resemble real administration and testing work. The focus is command composition, quoting, redirection, pipelines, and exit status---not complicated mathematics.
\end{labmode}

\medskip\noindent\textbf{Learning outcomes.}
\begin{enumerate}
  \item explain how a shell script composes existing programs;
  \item use variables, quoting, redirection, pipelines, and exit status safely;
  \item automate a repeatable OS observation instead of copying terminal commands manually.
\end{enumerate}

\medskip\noindent\textbf{Part A --- Generate a reproducible system report.}
\begin{lstlisting}[style=osbash]
cd experiments/Lab03
chmod +x *.sh
./system_report.sh
./system_report.sh my-report.txt
less my-report.txt
\end{lstlisting}
Identify which commands provide the kernel version, disk information, and process list. Explain why saving the report to a file makes the observation easier to compare later.

\medskip\noindent\textbf{Part B --- Query running processes.}
Try a process name that should exist, followed by one that should not:
\begin{lstlisting}[style=osbash]
./process_summary.sh bash
echo $?
./process_summary.sh definitely-not-a-process
echo $?
\end{lstlisting}
State what status 0 and a non-zero status mean in this script. Then start a harmless background process with \texttt{sleep 30 \&} and use the script to find it.

\medskip\noindent\textbf{Part C --- Summarise a log with standard tools.}
\begin{lstlisting}[style=osbash]
./log_summary.sh events.log
\end{lstlisting}
Trace how \texttt{grep}, \texttt{wc}, variables, and a loop cooperate. Add one new line to \texttt{events.log} and predict which count will change before running the script again.

\begin{keyidea}
Shell scripts are valuable in an OS course because they expose process creation, standard streams, exit status, and composition while letting existing programs perform most of the low-level work.
\end{keyidea}

\medskip\noindent\textbf{Part D --- Improve one tool.}
Choose one improvement: add a \texttt{--help} option, allow \texttt{system\_report.sh} to include memory information using \texttt{free -h}, or add an optional log level argument to \texttt{log\_summary.sh}. Demonstrate both a valid use and an invalid use.

\begin{verification}
Your submission should contain the report file, two process-query examples with exit statuses, the log summary before and after one controlled change, and a short note describing the script improvement you implemented.
\end{verification}

\begin{pitfall}
Quote variables that may contain spaces, for example \texttt{"\$file"}. Unquoted expansion can split one filename into several shell words and make a script behave unpredictably.
\end{pitfall}
\end{labproject}
